\section{Design Goals}
\label{sec:design-goals}

GAWorld v2 is intended as a configurable research instrument rather than as a
claim that a language model is, by itself, a valid population model.  Its
design goals therefore concern properties that a researcher can inspect in the
running system, and each property has a concrete mechanism in the present
implementation and a boundary that prevents the mechanism from being read as an
empirical guarantee.  Where v1 listed six properties, v2 lists eight, adding
two that the v1 platform could not satisfy: \emph{per-tenant world isolation}
and \emph{human participation under the same first-class intervention surface
as scheduled events}.

\subsection{Persistence across simulated time}
\label{sec:dgo:persistence}

An agent must carry consequences from one decision context into later ones,
rather than being reconstructed from a stateless prompt at every tick.  v2
combines a mutable state record with file-backed episodic memory, daily logs
and diaries, schedules, and longer-lived relationship and growth records.  A
stateful run restores the day counter and per-agent memories; the memory
lifecycle can retrieve, consolidate, summarize, and decay stored episodes.
Persistence is therefore testable: a researcher can inspect both the state
before and after a tick and the files used in later recall.  Persistence does
not establish psychological fidelity.  Many memories and summaries are
language-model outputs, configuration can disable lifecycle operations, and
several archived cross-phase evaluations are incomplete.

\subsection{Situatedness in a shared world}
\label{sec:dgo:situatedness}

Actions should be conditioned on where and when an agent acts and on what is
locally observable.  The simulation clock exposes day, time, and tick index;
the city map represents named locations, connections, and indoor rooms; and
environment, local-physical, social, policy, and information mechanisms
contribute context to the cognition loop.  Movement updates location before
later stages use the new setting, social influence is computed over the
current network rather than over an unconstrained cast of characters, and
indoor rooms are resolved against the same activity the agent is performing.
v2's addition of an \emph{away} state \citep{gaworldtravel} lets inhabitants
temporarily leave the city and return to a specific home node.  This is
situated computation, not a calibrated digital twin: map detail, physical
dynamics, agent access to information, and the indoor-room tree are
abstractions selected by configuration.

\subsection{Composability without concealing control flow}
\label{sec:dgo:composability}

Researchers should be able to change an experimental mechanism without
rewriting the complete agent loop.  v2 provides three bounded extension
surfaces: a configurable stage sequence for cognition, an event bus with
observe/contribute/filter semantics, and a registry for lifecycle-managed
plugins.  The audited built-in set now contains fourteen plugins: intervention,
skills, interests, life events, economy, local physical perception, real
work, dynamic behavior, spatial preferences, personality, indoor rooms and
venue capacity, organizations, governance, and offline citizens.  These join
existing memory, environment, and social-network subsystems that retain direct
compatibility paths.  Each plugin declares its dependencies, owns
configuration namespaces and per-agent state, and contributes through the bus
or through lifecycle hooks.  Composability is not unrestricted hot swapping:
the stage pipeline exchanges a shared step dictionary, and three legacy
subsystems (memory, environment, social network) retain direct paths.  The
audit table in \cref{sec:microkernel} lists the boundary state of each
subsystem.

\subsection{Controllability for counterfactual experiments}
\label{sec:dgo:controllability}

A research run should expose explicit points at which a scenario or
intervention is introduced.  Configuration fixes profiles, city assets,
random seed, mechanism flags, schedules, and policy events.  The
intervention plugin contributes a recommendation feed through the event bus
and records its domain-specific metrics.  The Controller API validates
structured movement requests through a priority-ordered validator chain and
exposes audited runtime interventions for state mutation, configuration
update, and agent removal.  v2 adds three new controllability surfaces: a
parallel-worlds engine that runs up to eight matched worlds with a shared
seed; a Research Workbench that compiles a research idea or a paper into a
pre-registered protocol with pre-specified outcomes and minimum-effect
thresholds; and a policy-simulation loop that proposes a candidate policy,
simulates, scores, and iterates.  None of these are guaranteed to identify a
real-world effect; they are designed to make the construct, the contrast, and
the completion status of every arm visible to the analyst.

\subsection{Inspectability from decision to artifact}
\label{sec:dgo:inspectability}

A claim about an artificial society should be traceable to artifacts below
the level of a prose summary.  v2 writes per-agent logs and memories,
environment timelines, state histories, plugin-owned outputs, and
experiment comparisons; its interfaces expose portions of those records for
inspection and replay.  The kernel Recorder additionally implements a
clock-stamped JSONL stream and records Controller interventions and denials.
v2 adds two new inspectability surfaces: a parallel-worlds results file that
records per-arm completion status, per-resident ATE, bootstrap confidence
intervals, and sign-flip $p$-values; and a per-tenant evidence ledger that
binds every reported number to a hash and an evidence class.  The unified
Recorder is not yet the main action/trace path: the established simulator
continues to write several specialized formats directly, and the
harness-built benchmark is a software-path diagnostic, not a replayable
trace.

\subsection{Backward-compatible modularization}
\label{sec:dgo:backcompat}

The system must remain usable while a monolithic research prototype is
decomposed.  A typed \texttt{Agent} view shares, rather than copies, the
legacy dictionary; \texttt{EventBus} preserves the legacy hook signature; and
the 12-stage pipeline was carved from the existing per-agent step while
retaining the shared step dictionary.  v2 adds two new compatibilities: the
\texttt{AccountStore} is optional (its absence reverts the system to single
mode), and the \texttt{DistributedRelayBackend} is optional (its absence
falls back to the in-process relay).  These adapters allow migrated and
legacy modules to coexist and make migration state observable.  They also
retain weakly typed surfaces and shared mutable state.  Backward
compatibility is a transition strategy, not evidence that modularization is
complete or that every old configuration produces an identical trajectory.

\subsection{Per-tenant world isolation}
\label{sec:dgo:multitenant}

A single deployment should host many independent societies, each with its
own city bundle, its own configuration overlay, its own run, its own output
tree, and its own visibility.  v2 implements this by treating a
\emph{World} as a first-class object: a city bundle plus a per-world
configuration overlay plus a per-world output root.  The active world is
resolved from a cookie, validated against the caller's role, and propagated
through a context variable.  All path-resolving functions
(\texttt{effective\_config}, \texttt{state\_csv\_path}, \texttt{records\_dir},
\ldots) consult the active world.  When the account store is absent, the
single \texttt{default} world is the only world and the system behaves as it
did before multi-tenancy was added.  A change in a world's output path is
therefore visible at the request level rather than at a process-restart
level.  Per-tenant isolation does not establish that two tenants running
identical configurations will produce identical trajectories; the underlying
model, the random seed, and the cross-process timing of distributed runs
remain sources of variation.

\subsection{Human participation as a first-class intervention}
\label{sec:dgo:participation}

Researchers and the public should be able to enter a society, claim a
resident, and act under the same first-class intervention surface that
scheduled events use.  v2 implements this through a \emph{player lease}:
browsers request a two-minute lease, see current perception, send a bounded
action, and release the resident.  In a multi-tenant deployment, the lease
is routed to the node that owns the resident; in a single-process
deployment, the lease is local.  The intervention type is the same, the
audit trail is the same, and the next-tick perception includes the player's
action.  This makes human participation inspectable in the same ledger as
scheduled events, but it does not validate that the human's action faithfully
represents the claimed resident.  The same evidence-class rules apply: a
player's claim is a real, inspectable intervention, not a ground truth.
